\section{Comprobaciones para la consolidación}\label{sec:comprobaciones}
En este anexo reúno las comprobaciones de las expresiones y mis criterios de
revisión. Distingo los resultados técnicos de mi experiencia al escribir y leer.

\subsection{Lectura de expresiones sin glosas visibles}
Incluyo en \path{lectura.tex} una hoja para practicar sin glosas. Contiene
las cuatro expresiones completas en otro orden, con una leyenda neutral y sin
traducciones resueltas. Se reconstruyó el alcance desde los operadores exteriores
y después se comparó con las lecturas del banco.

\begin{longtable}{@{}p{18mm}>{\raggedright\arraybackslash}p{\dimexpr\linewidth-18mm-2\tabcolsep\relax}@{}}
\toprule Caso & Lectura recuperada \\ \midrule\endhead
A / T05 & Existe un mismo libro que gusta a todas las personas del dominio:
\(\exists y\in C_2\;\forall x\in C_1:\ R_7(x,y)\). \\[4pt]
B / T02 & No son verdaderas ambas proposiciones:
\(\neg(P\land Q)\). La negación abarca la conjunción completa. \\[4pt]
C / T04 & A cada persona le gusta algún libro del dominio, que puede ser
distinto según la persona: \(\forall x\in C_1\;\exists y\in C_2:\ R_7(x,y)\). \\[4pt]
D / T03 & P es falsa y Q verdadera:
\((\neg P)\land Q\). La negación solo abarca P. \\[4pt]
\bottomrule
\end{longtable}
Las cuatro reconstrucciones coinciden con las claves. Los grupos exteriores
y anidados se distinguen en la hoja examinada. Tomo este resultado como una
revisión técnica con conocimiento previo del diseño y de las soluciones.
Todavía me queda evaluar mi lectura sin consultar las claves.

\subsection{Contrastes semánticos comprobados}
Se recalcularon las cuatro valoraciones booleanas mediante enumeración:
\begin{center}\begin{tabular}{cccc}
\toprule P & Q & T02 & T03 \\ \midrule
F & F & V & F \\
F & V & V & V \\
V & F & V & F \\
V & V & F & F \\
\bottomrule\end{tabular}\end{center}
En particular, con P y Q falsas, «no ambas» es verdadero, pero «no P y sí Q»
es falso. Esto identifica una diferencia de significado, no solo de dibujo.

También se enumeraron las 16 relaciones posibles entre dos personas y dos
libros: en 7 modelos ambas fórmulas son falsas, en 7 ambas son verdaderas y en
2 T04 es verdadera mientras T05 es falsa. Nunca ocurre el caso inverso.
Un contraejemplo al recíproco es que cada persona guste solo de un libro
distinto. Esta enumeración comprueba esos modelos finitos, no mide facilidad
de lectura ni constituye por sí sola una demostración general.

Se revisaron los dominios vacíos: sin personas ni libros, T04 es verdadera y
T05 falsa; sin personas pero con libros, ambas son verdaderas; con personas
pero sin libros, ambas son falsas. La gramática explicita ahora la lectura
de cuantificadores sobre dominios vacíos.

\subsection{Texto propio de varios enunciados: M01}
Para estudiar la continuidad del texto, planteo una inspección ficticia de
dos sensores en seis unidades distribuidas en dos páginas. Conservo el texto en
\path{datos/texto-integrado.tex}; la hoja sin soluciones está en
\path{relectura.tex}. Las dos presentaciones usan exactamente las mismas macros.

\begin{description}[style=nextline]
\item[U1] Según el registro de inspección, el sensor A está activo.
\item[U2] Según ese registro, el sensor B no está activo.
\item[U3] Según ese registro, A es rojo y B no es rojo.
\item[U4] Se define F como el conjunto de sensores rojos y activos en el
contexto fijado. La definición no afirma que exista alguno.
\item[U5] Se pregunta si existe un elemento de F que pertenezca al conjunto
de sensores. Como F se define mediante selección sobre C3, equivale a
preguntar si F es no vacío. La pregunta no afirma su respuesta.
\item[U6] La inspección termina antes de que comience el archivo:
\(12<13\), con tiempos medidos desde el mismo origen.
\end{description}

La revisión inmediata encontró los cinco referentes declarados y usados de
forma consistente, incluido el cambio de página. F se define antes de usarse
y las variables de U4 y U5 tienen ámbitos separados. Las negaciones reciben
proposiciones; las atribuciones y la pregunta permanecen en el nivel superior.
R8 recibe objeto y propiedad, R5 objeto y conjunto, y R1 dos eventos.

El texto hace visible una distinción necesaria: U1--U3 atribuyen información,
por lo que no autorizan concluir directamente \(F=\{A\}\). Esa igualdad sí
se sigue si, además, se acepta el contenido del registro como verdadero en
el contexto declarado. U6 se evalúa con los intervalos de la leyenda y no
afirma causalidad.

A partir de este ejemplo, precisé en la gramática la persistencia de referencias entre páginas,
la separación de unidades, el alcance de nombres definidos y la ausencia de
una secuencia de ejecución implícita entre párrafos. Estas son comprobaciones
estructurales inmediatas; no resultados de memoria al día siguiente.

\subsection{Mi próxima relectura}
Me queda volver a leer M01 otro día, sin consultar primero la explicación.
Usaré la hoja \path{relectura.tex} para anotar qué entiendo y después compararé
referencias, negaciones, definiciones y contexto. Así podré reconocer qué partes
recuerdo con facilidad y cuáles necesito ajustar.


\subsection{Consolidación editorial}
En el cuerpo del manual reúno los 22 signos, un orden orientativo de
trazado para cada uno, una gramática de referencia, los catálogos,
composición, M01, T01--T20 y la guía LaTeX. Los cambios gráficos, los registros
anteriores, los ejemplos iniciales y la investigación quedan en anexos.
Las gramáticas anteriores se conservan como fuentes históricas y no compiten
con la gramática vigente. El nombre de la escritura sigue siendo KalaScript.
